\section{项目的主要内容和技术路线}

\subsection{主要研究内容}

\par 本研究围绕基于智能体的拖延症干预系统展开，主要研究内容包括：

\par （1）智能体工作流架构设计：基于Spring AI Alibaba框架构建多智能体协作系统，采用图引擎进行工作流编排。系统配置多个推理型智能体，分别承担路由分发、简单/复杂项目对话、简单/复杂项目构建等职责，通过路由节点实现基于项目复杂度和聊天类型的智能路由。

\par （2）对话状态机与意图隔离机制：设计构建模式/对话模式/编辑模式三态交互模式，解决多轮对话中“任务构建意图”被闲聊打断的问题。构建模式中用户输入被视为需求补充，系统自动生成任务预览；对话模式允许基于已构建任务的自由探索；编辑模式支持任务结构的版本迭代。

\par （3）记忆持久化与上下文管理：采用基于 Redis 的断点持久化方案，实现基于Redisson分布式锁的多线程安全状态存储。配置摘要钩子进行长对话记忆摘要，当对话历史过长时自动触发摘要，保留最近若干条原始消息，确保上下文关联性与信息完整性。

\par （4）工具调用链可靠性保障：设计查询工具与操作工具分离架构，查询工具提供只读查询（获取项目、会话、标签、任务等），操作工具提供修改操作（创建/更新/删除任务等）。

\par （5）行为干预策略设计：基于微习惯理论和多巴胺设计原理，实现“方糖番茄钟”降低启动情绪成本；设计任务抽取随机推荐机制，利用变比率强化效应维持行为动机；构建积分与成就系统，通过即时反馈增强用户粘性。

\par （6）移动端前端设计：采用手绘风格营造温暖、亲切感，以黑白灰为主色调避免视觉压力。核心页面包括：首页（任务抽取入口）、任务清单页（紧急任务优先展示）、AI对话创建页（自然语言任务创建）、日历页（任务分布可视化）、专注页（方糖计时）、我的页（成就追踪）。

\subsection{技术路线}

\subsubsection{后端技术栈}

\par 后端采用Java 21 + Spring Boot 3架构，主要技术选型如下：

\begin{table}[htbp]
    \caption{后端主要技术选型}
    \begin{tabularx}{\linewidth}{c|X<{\centering}}
        \hline
        \textbf{技术领域} & \textbf{具体技术} \\ \hline
        编程语言 & Java 21 \\ \hline
        构建系统 & Maven 3 \\ \hline
        后端框架 & Spring Boot 3 \\ \hline
        数据库 & PostgreSQL 15 \\ \hline
        缓存 & Redis 7 \\ \hline
        ORM框架 & MyBatis Plus 3 \\ \hline
        安全认证 & Spring Security + JWT \\ \hline
        UUID生成 & UUID Creator \\ \hline
        AI框架 & Spring AI Alibaba \\ \hline
        对象映射 & MapStruct + ModelMapper \\ \hline
        号码识别 & libphonenumber \\ \hline
        短信验证 & dypnsapi20170525 \\ \hline
        智能文档 & Spring Doc \\ \hline
    \end{tabularx}
\end{table}

\subsubsection{智能体架构}

\par 智能体模块基于Spring AI Alibaba框架构建，核心架构层次包括：Controller层、Service层、AI引擎层、基础设施层、持久化层。

\par 工作流编排采用图结构，由路由节点到条件路由再到目标智能体。路由节点根据项目文本性质推断复杂度，根据项目元设置推断聊天模式/构建模式，返回对应的路由条件。

\subsubsection{数据持久化方案}

\par 采用PostgreSQL + Redis混合持久化架构：PostgreSQL存储用户、项目、任务、标签、会话、消息等核心业务数据，使用UUID主键和TIMESTAMPTZ时区；Redis存储活跃会话状态和检查点数据，支持多线程安全和自动TTL过期。

\par Session生命周期：用户首次发起对话时自动创建，每次对话完成后更新draftProjectView，支持分页查询和批量删除。Chat生命周期：对话开始时保存用户消息，AI生成过程中实时流式推送，对话完成后保存完整AI回复，新会话首次调用时从历史加载消息注入检查点。

\subsection{可行性分析}

\subsubsection{技术可行性}

\par Spring AI Alibaba框架提供了完整的智能体开发支持，包括推理型智能体、图引擎工作流编排、检查点持久化、摘要钩子等核心能力；Qwen和DeepSeek等国产开源大模型在中文理解和工具调用方面表现良好，能够满足任务拆解、对话交互等场景需求。

\par 阿里云提供了容器、数据库、号码认证服务、ECS等服务，从研发到测试到生产实现了全流程覆盖；PostgreSQL和Redis技术成熟稳定，社区支持完善。

\subsubsection{经济可行性}

\par 大模型API调用成本是主要运营支出。通过优化提示词设计、合理设置模型参数、采用流式输出减少等待时间等方式可有效控制成本。

\subsubsection{时间可行性}

\par 项目计划周期约12周，按“需求分析-架构设计-模块开发-系统集成测试”的逻辑主线推进。关键里程碑明确，各阶段任务可执行性强，能够在规定时间内完成系统开发和论文撰写。
